% -----------------------------------------------------------------------------------
\section{Inferring dependence from source}\label{sec:inference}

Theorem~\ref{thm:promiseopt} works because oracle substitution designates the positions at which the
agent is consulted. In an embedded environment, a closed program whose source may contain
predictors, simulations, or copies of the agent with no designated interface, the agent must
identify those positions itself. This section examines the three natural ways of doing so:
recognising the agent's code, recognising programs that compute the same function, and reasoning
about the agent's output in a formal theory. Each has a positive theorem on a restricted domain and a
characteristic failure outside it.

\subsection{Designated slots and literal copies}\label{sec:syntactic}

Suppose an environment is a template with designated agent slots, that syntactic replacement
changes exactly those slots, that every slot uses the same owned-query namespace, that all other
programs are constants, and that every finite-policy replacement terminates with an exact rational
payoff. Compiling the oracle calls of Theorem~\ref{thm:promiseopt} to these slots preserves every
counterfactual execution, so
\[
  U_{\mathrm{syn}}(E,T;A)=U_{\mathrm{oracle}}(E,A),
\]
and the uniform optimiser applies. This is a substantive positive answer: the rule evaluates
arbitrary code implementing the environment without analysing every consequence of its action; it
executes the stipulated replacements. It is also a restricted answer, since the representation
identifies the slots.

Without designated slots, literal recognition fails. Fix a globally total $T$ and let $g$ be a closed
program with $g(x)=T(x)$ for every $x$, encoded independently in a typed, atomic component so that its
representation contains no eligible occurrence of $T$'s code. (Wrapping $T$ is not enough, since the
original would remain an eligible nested occurrence.) Let the environment pay $\one[a\neq g(d(E))]$.
Literal substitution changes the live agent and leaves $g$ unchanged, so $T$ receives $0$ while an
agent taking the opposite action receives $1$. Every computation terminates. This is the fixed-copy
diagonal of Theorem~\ref{thm:fixedcopy} again, now evading a syntactic matcher. Literal matching also
violates $\cls_{\mathrm{ext}}$: two extensionally equal reference agents with different source code
produce different replacement masks. A slot-based representation avoids the problem by preserving
the agent's role independently of the code that currently occupies it; code equality and role
identity should not be conflated.

\subsection{Extensional matching}\label{sec:extensional}

The natural repair is to replace every component that computes the same function as the agent.
To make this precise, let $E$ contain a finite list of typed, atomic components $g_j$, fix comparison
domains $X_j$ and a reference function $t$, and let
\[
  S_t(E)=\{j: g_j|_{X_j}=t|_{X_j}\}.
\]
To evaluate $A$, replace exactly the components in $S_t(E)$ by $A$ and execute; write
$U_{\mathrm{ext}}(E,t;A)$ for the result. Matching occurs once, before replacement; component
boundaries and the description supplied to the agent stay fixed; and every permitted counterfactual
execution is required to terminate, which termination of the original environment does not imply.
The principal choices are \emph{local} equality, on a specified finite set or a frozen baseline
trace, and \emph{global} equality, on the entire input domain.

\begin{proposition}[Matching on the counterfactual trace]\label{prop:tracematch}
Requiring replacement exactly when a component agrees with $t$ on the inputs used in the resulting
counterfactual execution can leave no consistent replacement set, although every candidate
execution terminates.
\end{proposition}

\begin{proof}
Let $t(0)=t(1)=0$, $g(0)=0$, $g(1)=1$, and $A(0)=1$, and let the environment call $g(0)$ and then call
$g(1)$ exactly when the first result is $1$. Without replacement only the input $0$ is used, where $g$
agrees with $t$, so replacement is required. With replacement, $A(0)=1$ triggers the call on $1$, where
$g$ disagrees with $t$, so replacement is forbidden.
\end{proof}

\enquote{The inputs it is run on} must therefore name the execution that determines those inputs.
Equality of total programs on a fixed finite set is decidable by running them, so local matching
needs no appeal to Rice's theorem. It fails for a different reason.

\begin{theorem}[Local equality can leave no optimal agent]\label{thm:localmatch}
Under frozen matching on the single input $d(E)$, there is a terminating environment satisfying
$\cls_{\mathrm{self}}$ and $\cls_{\mathrm{ext}}$ on which no deterministic agent satisfies anchored
optimality~\eqref{eq:anchored}.
\end{theorem}

\begin{proof}
Let $E$ contain the components $g_0\equiv0$ and $g_1\equiv1$, call both at $d(E)$, query the agent
there, and pay $1$ exactly when the two component outputs agree. For a reference agent with
$t(d(E))=a$, exactly $g_a$ matches. With $t$ itself substituted, the outputs are $a$ and $1-a$, so
$U_{\mathrm{ext}}(E,t;t)=0$. With an agent $A$ taking $1-a$, the matched component $g_a$ becomes $A$
and both outputs equal $1-a$, so $U_{\mathrm{ext}}(E,t;A)=1$.
\end{proof}

This is not an undecidability theorem. It is a pathology of identifying computations through
accidental agreement on one input, and it should not be presented as a theorem about functional
decision theory without that identification.

\begin{proposition}[Rematching does not compose]\label{prop:rematch}
Let $R_{t\to a}$ replace by $a$ the components that match $t$. With components $g_0\equiv0$,
$g_1\equiv1$ and reference $t\equiv0$, and with matches recomputed after each replacement,
$R_{a\to b}(R_{t\to a}(E))\neq R_{t\to b}(E)$ for $a\equiv1$ and $b\equiv0$.
\end{proposition}

\begin{proof}
$R_{t\to a}$ turns $g_0$ into $1$, giving the pair $(1,1)$. Both components now match $a$, and $R_{a\to b}$
gives $(0,0)$. Direct replacement of $t$ by $b=t$ leaves $(0,1)$.
\end{proof}

Keeping an original dependency label attached to the first component restores composition, but the
label is additional structure; equality of the current functions does not recover it. Global
matching avoids accidental agreement, but it identifies the wrong thing.

\begin{theorem}[Global equality leaves the diagonal intact]\label{thm:complement}
For globally total agents, global-equivalence replacement of atomic components admits no universal
optimiser when fixed algorithms are allowed.
\end{theorem}

\begin{proof}
Fix $T$, compile an atomic program $h_T(x)=1-T(x)$, and let the environment pay
$\one[a=h_T(d(E))]$. Since $h_T$ disagrees with $T$ everywhere, it is not replaced. $T$ receives $0$,
an agent taking the opposite action receives $1$, and all calls terminate.
\end{proof}

The diagnosis is that \enquote{depends functionally on $T$} is not \enquote{is extensionally equal
to $T$}. If the intended counterfactual turns $1-T$ into $1-A$, it transports dependence through a
program context, which is a stronger operation than equality-based replacement. Global matching
also brings in undecidability, and the expected route needs one correction. For a uniformly total
family $g_e$ and a fixed total $T$, non-equivalence has a finite witness, so $\{e:g_e=T\}$ is co-c.e.;
one cannot uniformly make equality hold exactly when a machine halts, but one can make it hold exactly
when it does not.

\begin{theorem}[Global matching can decide halting]\label{thm:haltred}
Assume atomic global matching, frozen descriptions, and globally total binary agents, and let $P$ be a
decidable class of descriptions that contains all the environments constructed below and omits
arbitrarily long strings. No agent satisfies anchored optimality on all environments in $P$.
\end{theorem}

\begin{proof}
Fix a candidate $T$ and let
\[
  g_e(x)=
  \begin{cases}
    1-T(x), & x\notin P\text{ and machine }e\text{ halts within }|x|\text{ steps},\\
    T(x),   & \text{otherwise}.
  \end{cases}
\]
Each $g_e$ is total, and $g_e=T$ exactly when $e$ does not halt. Let $E_e$, with $d_e=d(E_e)\in P$,
receive the action $a$, compute $b=g_e(d_e)$, and pay $u(a,b)=2(1-b)+a$:
\begin{center}
\small
\begin{tabular}{@{}lcc@{}}
  \toprule
  & action $0$ & action $1$ \\
  \midrule
  $b=0$ & 2 & 3 \\
  $b=1$ & 0 & 1 \\
  \bottomrule
\end{tabular}
\end{center}
Because $d_e\in P$, $g_e(d_e)=T(d_e)$ whether or not $e$ halts. If $e$ does not halt, global matching
replaces $g_e$ by the candidate $A$, so $b=A(d_e)$ and the values of the two actions are $2$ and $1$;
optimality requires $T(d_e)=0$. If $e$ halts, $g_e$ stays fixed, and with $b$ held fixed action $1$
pays one unit more; optimality requires $T(d_e)=1$. The computable map $e\mapsto T(d_e)$ would decide
the halting problem.
\end{proof}

The payoff table is Newcomb's, and the reduction works by making it undecidable whether its
predictor counts as linked to the agent. Every execution terminates and every payoff is rational,
but the result uses global equality, atomic components compiled from $T$, and a class already
admitting fixed-constant diagonals; it satisfies executed-query locality and violates kernel
locality. It is a valid reduction, not a semantics-independent impossibility arising from the mere
presence of embedded programs. Under local matching $g_e$ always matches at $d_e$, so action $0$ always
wins and the reduction disappears. Two further qualifications bound any positive subclass. Clocks do
not decide global equivalence, since even a program that simulates machine $e$ for only $|x|$ steps
can equal zero everywhere exactly when $e$ never halts. And \enquote{replace when a theory $S$ proves
equality} changes the semantics: it agrees with global matching only where certified matching
information is complete, since positive certificates do not show that uncertified components are
unequal.

\subsection{Proofs and logical counterfactuals}\label{sec:logical}

Three questions must be separated: whether proofs can establish the values of independently
specified counterfactual programs; whether implications about the actual agent can serve as
counterfactual values; and whether logical uncertainty supplies useful finite-time estimates. The
answers are yes, no, and not by itself.

\begin{proposition}[Ground proof-based optimisation]\label{prop:proofopt}
Let $S$ be a sound, effectively axiomatised arithmetic theory that verifies finite program traces.
The agent that, on an environment in $\mathcal P_{\mathbb Q}$, enumerates $S$-proofs until it has, for
every $\pi\in\Pol_E$, a proof of $\mathrm{Eval}(E,\pi)=r_\pi$ for some rational $r_\pi$, and then
maximises the $r_\pi$, is promise-total and optimal under $O_{\mathrm{oracle}}$.
\end{proposition}

\begin{proof}
Each terminating execution $E^\pi$ has a finite trace, the trace has a proof in $S$, and soundness
makes every proved value correct. Only finitely many proofs are needed.
\end{proof}

The incompleteness of a fixed theory does not prevent it from verifying every particular terminating
computation; what it may fail to prove is a uniform termination statement or a global equivalence.
The proposition evaluates the same counterfactual programs as Theorem~\ref{thm:promiseopt} and infers
no new semantics from proof search.

Material implication is different. Write $H_a$ for $T(d(E))=a$ and $Y_u$ for $U(E,T)=u$. If
$S\vdash\neg H_a$, then $S\vdash H_a\to Y_u$ for every $u$, so the relation
\[
  \mathcal V(E,T,a)=\{u: (H_a\to Y_u)\text{ is true}\}
\]
contains only the actual payoff at the actual action and every rational at every other action. It is
a relation, not an operator; any numerical counterfactual drawn from it needs a selection principle
that the implication does not supply. For an actually terminating computable agent, a sufficiently
strong theory verifies the agent's execution and so refutes every non-actual antecedent, and
conditioning a coherent probability over its models on such an antecedent is undefined. Proof search
can also select a spurious value, as in the \enquote{five-and-ten} problem
\autocite[§2.1]{demskigarrabrant2019}, which can be realised by a total program.

\begin{theorem}[A computable five-and-ten failure]\label{thm:fiveten}
Let $E_0$ pay the action chosen from $\{5,10\}$, and let $S$ be an effectively axiomatised theory
satisfying the derivability conditions for Löb's theorem \autocite{lob1955}. There is a globally
total agent $T$ that, on $d(E_0)$, takes $5$ because it has found a proof that taking $10$ yields $0$.
It receives $5$, while the constant agent $10$ receives $10$.
\end{theorem}

\begin{proof}
By the recursion theorem let $T$ return $10$ on descriptions other than $d(E_0)$, and on $d(E_0)$
enumerate $S$-proofs until it finds a proof of
\[
  \psi=\bigl(T(d(E_0))=5\to U(E_0,T)=5\bigr)\wedge\bigl(T(d(E_0))=10\to U(E_0,T)=0\bigr),
\]
and then return $5$. Let $q$ be the sentence $T(d(E_0))=5$. The definitions of $T$ and $E_0$ give
$S\vdash q\to\psi$, since under $q$ the first conjunct holds and the second is vacuous. Hence
$S\vdash\Prov_Sq\to\Prov_S\psi$. Formalising the search gives $S\vdash\Prov_S\psi\to q$. So
$S\vdash\Prov_Sq\to q$, and Löb's theorem yields $S\vdash q$, whence $S\vdash\psi$. A proof of $\psi$
therefore exists, the search halts, and $T$ returns $5$.
\end{proof}

This refutes the displayed implication-based selection scheme, not every logical-counterfactual
procedure. \textcite[§2]{bensontilsen2014} analyses the mechanism, and his classification of
outcomes by proof-search order concerns a specified transparent-universe model with a provability
oracle; it should not be quoted as a theorem about every computable proof-based agent. A different
response, the \enquote{chicken rule}, has the agent take any action it can prove it will not take
\autocite{slepnev2011}; within Benson-Tilsen's model it is the only universally successful search
strategy \autocite[Corollary~4.6]{bensontilsen2014}, and it blocks spurious proofs at the cost of
counterfactuals that can differ from what another agent would actually obtain
\autocite{fallenstein2015chicken}.

A precise logical semantics that avoids spurious values does so by becoming partial.

\begin{definition}[Guarded modal counterfactuals]\label{def:guarded}
Let $(A_i)$ and $(U_j)$ be closed modal formulas, provably mutually exclusive and exhaustive, describing
the actions and outcomes of an agent and a problem. Action $i$ \emph{leads to} outcome $u_j$ if
\[
  \mathrm{GL}\vdash A_i\to U_j\qquad\text{and}\qquad\mathrm{GL}\nvdash\neg A_i,
\]
and has no value otherwise \autocite{fallenstein2015third}.
\end{definition}

The outcome is unique when defined and is undefined if the action is provably impossible or no
outcome is determined; the semantics \enquote{fails visibly} rather than assigning a spurious value.
Fallenstein uses it to state the optimality of modal updateless decision theory on problems where
every action has a value \autocite{fallenstein2015third}. This is optimality relative to the specified
counterfactuals, not relative to physical agent replacement, and the two can diverge: on a variant of
five-and-ten where every action has a value, these counterfactuals certify as optimal an action that
another agent would beat \autocite{fallenstein2015chicken}. In provability logic $\mathrm{GL}$
\autocite{boolos1993} the guard is informative. For ordinary programs it is not.

\begin{proposition}[Guarded counterfactuals for terminating programs]\label{prop:guardedarith}
Replace $\mathrm{GL}$ by an arithmetic theory $S$ that verifies terminating computations and proves the
determinism of its interpreter. For an agent that terminates on $d$, every non-actual action has no
value.
\end{proposition}

\begin{proof}
If the actual action is $b$, the finite execution trace gives $S\vdash T(d)=b$, and determinism gives
$S\vdash T(d)\neq a$ for every $a\ne b$, so every non-actual action fails the guard.
\end{proof}

This reconstruction is not attributed to the source of Definition~\ref{def:guarded}; it explains
why positive results for modal agents, or agents with provability oracles, do not transfer directly
to ordinary terminating programs.

The modal setting has its own diagonal and its own optimality theorem, and their quantifiers matter.
Consider fixed, fully modalised decision schemata $T(\mathbf u)$ and problems that are \emph{provably
extensional}, in the sense that $\mathrm{GL}\vdash(\mathbf a\leftrightarrow\mathbf b)\to
(P(\mathbf a)\leftrightarrow P(\mathbf b))$. Every such fixed theory with at least two actions and
two outcomes fails on some provably extensional problem \autocite{fallenstein2014evil}: the problem
rewards any action other than the one the theory takes when facing the constant second-best outcome,
modal fixed-point uniqueness makes that action the theory's own, and a different constant action
wins. This is the fixed-constant diagonal in modal form. Conversely, for every finite provably
extensional problem $P$ there is a true sentence $\varphi$ such that modal updateless decision theory
using $\mathrm{PA}+\varphi$ selects an optimal action on $P$
\autocite{fallenstein2014optimality,fallenstein2014improving}. The quantifiers are
$\forall P\,\exists\varphi$: there is no single proof system that works for all problems, and the
general procedure of \textcite{soaresfallenstein2015} that queries proofs uses a halting oracle. $\mathrm{GL}$
itself is decidable, so none of this shows that restricted modal computation is uncomputable.

Logical induction \autocite{garrabrant2016li} supplies computable, improving assessments of logical
sentences. It does not specify which counterfactual program or intervention to evaluate, and it
imposes no exact constraint on any finite stage: by its Theorem~4.6.1, changing finitely many market
days preserves the logical induction criterion. A decision procedure can use an inductor to estimate
the values of policy-substituted programs, but Theorem~\ref{thm:exactsel} still applies to any
resulting computable rule. Decision procedures built on logical induction typically add
\emph{exploration}, taking an action whenever its probability falls below a threshold so that the
conditionals on it remain defined \autocite{garrabrant2017obstacles}; Section~\ref{sec:candidates}
shows that what such a procedure computes depends on the exploration scheme it is given.

\subsection{What inference supplies}

The pattern across these results is uniform. Wherever a substitution is designated---by an oracle
interface, by slots, or by an evaluator whose results a proof certifies---optimisation is
computable on the corresponding termination promise. Wherever the agent must decide for itself
which parts of a closed program would change with its output, each recognition rule fails in its
own way: literal matching misses encoded copies, local matching manufactures dependence from
accidental agreement, global matching misses transformed dependence and can be undecidable, and
implication-based reasoning about the actual agent yields either spurious values or none. None of
these failures is a theorem that no coherent rule exists. They show that the rule is the open
problem, which the next section takes up axiomatically.
